\documentclass[11pt,letterpaper]{article}

\usepackage[margin=1in]{geometry}
\usepackage[T1]{fontenc}
\usepackage[utf8]{inputenc}
\usepackage{lmodern}
\usepackage{amsmath}
\usepackage{amssymb}
\usepackage{booktabs}
\usepackage{enumitem}
\usepackage[hidelinks]{hyperref}
\usepackage{setspace}

\setstretch{1.08}
\setlength{\parindent}{0pt}
\setlength{\parskip}{5pt}

\title{ECE 280 Lab 3 -- TA Solution and Troubleshooting Manual\\Signals and Systems: CT/DT System Properties Explorer}
\author{}
\date{}

\begin{document}
\maketitle

Use this document as a fast grading and troubleshooting guide for:
\begin{itemize}[leftmargin=1.2em]
    \item Lab 3 manual and report template submissions
    \item System-property verification from hardware-acquired data
    \item TA checkoffs during in-lab demonstrations
\end{itemize}

\section{Quick Pass/Flag Menu}

\subsection*{Pass Immediately If All Are True}
\begin{itemize}[leftmargin=1.2em]
    \item Student presents physical evidence for all three circuits (RC, diode clipper, comparator).
    \item All required screenshots include student name/ID and are legible.
    \item MATLAB verification uses measured input/output data (not synthetic-only arrays).
    \item Superposition and TI tests include clear overlays and reported NRMSE values.
    \item Final truth table is evidence-backed and consistent with measured behavior.
\end{itemize}

\subsection*{Flag for Review If Any Are True}
\begin{itemize}[leftmargin=1.2em]
    \item Property claims are made without scope evidence and/or numerical checks.
    \item No unique test inputs (students reuse common canned waveforms without assigned seed variation).
    \item Data appears copied across circuits (identical traces for RC/clipper/comparator).
    \item MATLAB script cannot identify hardware data file origin.
\end{itemize}

\section{Instructor Key: Expected Property Outcomes}

Assume ideal operating conditions and correctly built circuits.

\begin{table}[h!]
\centering
\caption{Reference Truth Table (Expected in Most Correct Builds)}
\begin{tabular}{lcccc}
\toprule
System & Linear & Time-Invariant & Causal & BIBO Stable \\
\midrule
RC low-pass filter & Yes & Yes & Yes & Yes \\
Diode clipper & No & Yes & Yes & Yes \\
Comparator (fixed threshold) & No & Yes & Yes & Yes \\
\bottomrule
\end{tabular}
\end{table}

\textbf{Notes for TAs:}
\begin{itemize}[leftmargin=1.2em]
    \item Clipper and comparator are memoryless nonlinear maps, so they generally remain TI but violate superposition.
    \item Comparator with drifting threshold or supply instability may appear weakly non-TI in practice; this should be discussed, not ignored.
\end{itemize}

\section{Pre-Lab Solution Guidance}

\subsection*{Required Conceptual Targets}
\begin{itemize}[leftmargin=1.2em]
    \item Correct formal definitions of linearity, TI, causality, and BIBO stability.
    \item Predicted outcomes for all three circuits with mechanism-based justification.
    \item Test design includes two distinct inputs, nontrivial coefficients $a,b$, and a measurable time shift $\Delta$.
\end{itemize}

\subsection*{Accept/Revise Criteria}
\begin{itemize}[leftmargin=1.2em]
    \item Accept if definitions are mathematically valid and predictions are physically coherent.
    \item Revise if student confuses ``bounded output'' with ``low noise'' or treats TI as ``periodic output.''
\end{itemize}

\section{Measurement Targets and Acceptance Bands}

\subsection*{RC Filter Baseline}
\begin{itemize}[leftmargin=1.2em]
    \item Expect attenuation and phase lag dependent on frequency relative to cutoff.
    \item Accept if gain/phase trends are physically correct even with moderate component tolerance.
\end{itemize}

\subsection*{Diode Clipper Baseline}
\begin{itemize}[leftmargin=1.2em]
    \item Clipping threshold typically near diode conduction region (depends on topology/bias).
    \item Accept if waveform clearly shows limited peaks and threshold annotation is reasonable.
\end{itemize}

\subsection*{Comparator Baseline}
\begin{itemize}[leftmargin=1.2em]
    \item Output should switch between rails (or logic levels) around threshold crossing.
    \item Accept if transition behavior is consistent with threshold and input crossing events.
\end{itemize}

\section{Numerical Verification Key (MATLAB)}

\subsection*{Linearity Test Expectations}
Students test:
\[
\mathcal{T}\{a x_1 + b x_2\}\quad\text{vs}\quad a\mathcal{T}\{x_1\}+b\mathcal{T}\{x_2\}
\]

Expected trends:
\begin{itemize}[leftmargin=1.2em]
    \item RC filter: low residual, small NRMSE (often \(<10\%\) with good alignment).
    \item Diode clipper: clear failure where clipping activates; moderate/high NRMSE expected.
    \item Comparator: strong failure (binary nonlinearity), usually high NRMSE except trivial small-signal regimes.
\end{itemize}

\subsection*{Time-Invariance Test Expectations}
Students test:
\[
\mathcal{T}\{x(t-\Delta)\}\quad\text{vs}\quad y(t-\Delta)
\]

Expected trends:
\begin{itemize}[leftmargin=1.2em]
    \item All three systems should be TI in a stable setup with fixed parameters.
    \item Elevated TI NRMSE often indicates poor synchronization, trigger drift, or threshold drift rather than true TI failure.
\end{itemize}

\section{TA Checkoff Script (2--3 Minutes per Team)}
\begin{enumerate}[leftmargin=1.4em]
    \item Ask team to state predicted truth table before showing data.
    \item Verify one scope capture per circuit with visible settings and student ID.
    \item Ask for one MATLAB linearity overlay and one TI overlay (any circuit).
    \item Ask: ``Why did this property pass/fail physically?''
    \item Confirm final report will include handwritten truth table and code/data linkage.
\end{enumerate}

\section{Troubleshooting Matrix}

\begin{table}[h!]
\centering
\caption{Common Failure Modes and Corrective Actions}
\begin{tabular}{p{0.24\textwidth} p{0.33\textwidth} p{0.35\textwidth}}
\toprule
\textbf{Observed Symptom} & \textbf{Likely Cause} & \textbf{TA Action} \\
\midrule
No clipping visible in clipper output & Diode reversed, incorrect reference node, or amplitude too low & Check diode orientation and threshold path; increase input amplitude above expected clipping level. \\
Comparator output stuck high/low & Threshold out of range, missing ground reference, wiring error & Verify common ground and reference divider; sweep threshold and confirm crossing events. \\
RC output looks identical to input at all frequencies & Cutoff too high or probe connected before RC node & Recheck node location and component values; test with higher input frequency to reveal filtering. \\
TI test fails for all circuits & Time alignment/trigger mismatch, not intrinsic circuit behavior & Enforce common trigger method and post-alignment in MATLAB before property claim. \\
Linearity appears to pass clipper/comparator & Inputs too small to engage nonlinearity & Increase amplitude to activate clipping/switching region and repeat superposition test. \\
MATLAB NRMSE unstable/run-to-run variation & Inconsistent sample window, serial timing jitter, poor preprocessing & Use fixed acquisition length, consistent trigger point, detrend/align before computing metrics. \\
\bottomrule
\end{tabular}
\end{table}

\section{Grading Menu Aligned to Lab 3 Rubric (100 pts)}
\begin{itemize}[leftmargin=1.2em]
    \item Pre-lab quality (15): definitions, predictions, and test design are rigorous.
    \item Circuit implementation/evidence (20): all three circuits built and documented.
    \item Linearity verification (15): valid superposition workflow and justified conclusion.
    \item TI verification (15): proper shifted-input methodology and justified conclusion.
    \item Causality/stability evaluation (15): bounded tests and anticipatory behavior check.
    \item Truth table/technical reasoning (10): final table is coherent and evidence-backed.
    \item Code/data authenticity (10): MATLAB clearly uses measured hardware data.
\end{itemize}

\subsection*{Suggested Deductions}
\begin{itemize}[leftmargin=1.2em]
    \item Missing ID on required evidence: \(-3\) to \(-10\), depending on extent.
    \item No quantitative metric (NRMSE or equivalent): \(-4\) to \(-8\).
    \item Property conclusions unsupported by data: \(-5\) to \(-12\).
    \item Synthetic-only MATLAB verification: \(-8\) to \(-15\).
\end{itemize}

\section{Minimum Passing Submission Standard}
A minimally acceptable submission should include:
\begin{itemize}[leftmargin=1.2em]
    \item Physical evidence of all three circuits.
    \item At least one valid linearity and one valid TI test workflow.
    \item Final truth table with short justifications per system.
    \item MATLAB script excerpts tied to acquired datasets.
    \item Student identification visible in required screenshots.
\end{itemize}

\end{document}
